% ============================================================
% Capitolo 1 — Introduzione (traduzione italiana di
% chapters/ch01_introduction.tex; label invariati)
% ============================================================
\chapter{Introduzione}
\label{ch:intro}

\section{Motivazione}
\label{sec:motivation}

L'apprendimento per rinforzo (\emph{reinforcement learning}, RL) addestra
politiche di controllo per tentativi ed errori e, in qualunque dominio non
banale, il progettista affronta una questione di scheduling prima ancora
del primo passo di gradiente: quando i task da padroneggiare coprono un
intervallo di difficoltà, \emph{in quale ordine vanno presentati}? Due
risposte sono naturali. Il \emph{curriculum learning} presenta i task dal
facile al difficile, sulla falsariga della pedagogia umana: costruire
competenza dove il successo è raggiungibile, poi estenderla.
L'\emph{addestramento misto} presenta tutto fin dall'inizio --- l'analogo
diretto dell'addestramento su un dataset mescolato, che è il default
indiscusso nell'apprendimento supervisionato. Entrambe le strategie
consumano lo stesso budget; scegliere tra le due è, in linea di principio,
gratis. Se la scelta cambia ciò che la policy appresa fa, essa è tra le
leve più economiche a disposizione del progettista --- e tra le meno
caratterizzate, specialmente nei contesti multi-agente, dove la maggior
parte delle assunzioni della letteratura sui curricula (apprendista
singolo, reward sparse, ambiente stazionario) non vale.

La guida autonoma urbana è un banco di prova esigente per porre la
domanda. Offre un asse di difficoltà naturale (la lunghezza del
percorso), agenti eterogenei con embodiment e vincoli diversi (veicoli e
pedoni), una tensione intrinseca tra prestazione e sicurezza (policy più
veloci completano più percorsi e urtano più cose) e un test
out-of-distribution canonico (guidare su una città mai vista in
addestramento). È inoltre intrinsecamente multi-agente: le policy non si
limitano a condividere un mondo, ma plasmano l'ambiente l'una dell'altra
mentre apprendono --- la non-stazionarietà al cuore dell'apprendimento
per rinforzo multi-agente. Questa tesi pone lì la domanda
sull'ordinamento, e tratta il regime di addestramento stesso come
variabile sperimentale.

\section{Domanda di ricerca e ipotesi}
\label{sec:rq}

\begin{quote}
  \emph{Il curriculum learning da facile a medio a difficile produce un
  comportamento di guida multi-agente misurabilmente diverso rispetto
  all'addestramento mixed-batch, a parità di budget di addestramento?}
\end{quote}

La domanda riguarda deliberatamente il \emph{comportamento}, non una
singola classifica di tassi di successo: il confronto è valutato lungo
quattro assi --- efficienza campionaria, prestazione finale con la sua
composizione dei fallimenti, dinamica di addestramento entro la run e
generalizzazione di mappa --- sempre separatamente per veicoli e pedoni.
Tre ipotesi falsificabili strutturano l'analisi:

\begin{itemize}
  \item \textbf{H1 (curriculum).} L'esposizione graduale produce una
        migliore efficienza campionaria iniziale --- i task facili
        offrono successi raggiungibili da cui partire --- con il rischio
        noto di plateau sui task difficili introdotti tardi, o del loro
        oblio.
  \item \textbf{H2 (misto).} L'esposizione uniforme copre l'intera
        distribuzione dei task dal passo zero, al costo di un segnale
        iniziale più rumoroso dominato dai fallimenti sui task
        difficili; il suo payoff atteso è la robustezza --- ciò che si
        vede in addestramento non cambia nel tempo.
  \item \textbf{H3 (interazione, secondaria).} L'entità --- forse la
        direzione --- dell'effetto di regime può dipendere da come viene
        stimato il valore. Nelle architetture con addestramento
        centralizzato il critic è la componente che digerisce la
        struttura multi-agente; sostituirne l'encoder (MLP piatto vs.\
        message passing su grafo tra agenti) può cambiare ciò che
        l'ordinamento dei task produce.
\end{itemize}

Le risposte in breve della campagna, sviluppate nel \cref{ch:results}:
la clausola di efficienza di H1 fallisce (tutte le celle apprendono
ugualmente in fretta all'inizio --- il reward shaping denso fornisce già
il segnale che i curricula di solito procurano) mentre l'ordinamento
seleziona invece \emph{quale} policy emerge, cauta oppure veloce e
incline alle collisioni; la clausola di robustezza di H2 fallisce (la
copertura più ampia non ha comprato alcun vantaggio di trasferimento); e
H3 regge in forma forte --- con un critic GNN l'effetto di regime
in-domain svanisce del tutto, e solo il vantaggio di generalizzazione
del curriculum sopravvive in entrambe le architetture.

\section{L'approccio in sintesi}
\label{sec:approach}

Tutti gli esperimenti girano in CARLA 0.9.16 sulla mappa urbana Town03,
con Town05 riservata al test di generalizzazione. Sei agenti apprendono
simultaneamente in un unico mondo condiviso: tre veicoli e tre pedoni,
ciascuno con un proprio percorso assegnato a ogni episodio --- un design
che chiamiamo \emph{copertura parallela dei task} (parallel task
coverage), che moltiplica l'esperienza per episodio e rende
l'interazione tra agenti un fenomeno di prim'ordine anziché un
disturbo. È la cornice fissa di entrambi i bracci sperimentali, non una
variabile. Le policy sono addestrate con MAPPO in regime di
addestramento centralizzato con esecuzione decentralizzata: una policy
condivisa per classe e un critic centralizzato che vede le osservazioni
di tutti gli agenti solo durante l'addestramento. La difficoltà è la
lunghezza target del percorso (30/60/\SI{100}{m}); il successo è il
completamento rigoroso del percorso.

I due regimi condividono ogni componente tranne la regola di selezione
del livello. Il braccio \emph{curriculum} parte dal facile e sblocca i
livelli più difficili attraverso un gate di competenza a finestra (con
cap di pressione basati sul budget che garantiscono la progressione); il
braccio \emph{batch} alterna uniformemente tutti i livelli fin dal primo
episodio. La campagna finale è un fattoriale $2 \times 2$ --- \{critic
MLP, GNN\} $\times$ \{curriculum, batch\} --- di quattro run da 3 milioni
di passi con un unico seed, appaiate tra le celle, ognuna seguita da una
valutazione di 400 episodi che include Town05. Con una run per cella,
ogni contrasto è un caso di studio appaiato per seed: gli effetti sono
letti contro una banda di varianza a seed identico misurata su repliche
identiche, e il contrasto primario è osservato due volte, una per
architettura del critic. La campagna poggia su una più lunga fase di
sviluppo guidata da gate --- documentata per intero, inclusi i suoi
fallimenti e i bug d'ambiente che ha scoperto --- che è parte del
contenuto metodologico della tesi (\cref{ch:methodology}).

\section{Contributi}
\label{sec:contributions}

\begin{enumerate}
  \item \textbf{Un framework di confronto controllato a parità di
        budget} per i regimi di addestramento nella guida multi-agente:
        run appaiate per seed con configurazioni byte-identiche a meno
        dei flag sperimentali, regole di misura rigorose per agente e
        una valutazione a tre livelli (flusso di addestramento, ricalcolo
        statico, valutazione finale del checkpoint con trasferimento di
        mappa).
  \item \textbf{Una metodologia di sviluppo guidata da gate con un
        registro onesto}: candidati a manopola singola promossi o
        revertiti contro gate quantitativi predefiniti, meccanismo ed
        esito registrati separatamente, e difetti di validità
        dell'ambiente --- incluso uno che ha spostato la metrica
        principale di $\approx 50$ punti percentuali senza alcun
        cambiamento di policy --- documentati con le loro firme di
        rilevamento.
  \item \textbf{Una caratterizzazione empirica di che cosa fa
        l'ordinamento dei task} in questo contesto: nessun effetto di
        efficienza iniziale, ma selezione tra equilibri di guida
        qualitativamente diversi; un vantaggio del curriculum sul
        criterio di successo rigoroso con il critic MLP; pedoni
        insensibili al regime; e un'erosione tardiva di piattaforma
        misurata a livello di task fissato.
  \item \textbf{Evidenza che gli effetti di regime sono condizionati
        dalla stima del valore} (H3): sostituire l'encoder del critic da
        MLP a GNN cancella il vantaggio in-domain del curriculum,
        recupera la cella batch e lascia la generalizzazione di mappa
        come unica componente che si replica tra le architetture.
\end{enumerate}

\section{Struttura della tesi}
\label{sec:outline}

Il \cref{ch:background} introduce il background necessario:
apprendimento per rinforzo e PPO, apprendimento multi-agente con critic
centralizzati, curriculum learning e il suo controllo ad addestramento
misto, e il simulatore CARLA. Il \cref{ch:system} descrive la
piattaforma sperimentale --- ambiente, spazi di osservazione e azione,
design della reward, generazione dei percorsi, i due regimi di
addestramento e le varianti del critic. Il \cref{ch:methodology}
definisce le regole di misura e il protocollo guidato da gate, documenta
gli episodi di validità dell'ambiente e specifica il design della
campagna. Il \cref{ch:results} riporta i risultati lungo i quattro assi
di confronto, chiudendo con i verdetti su H1--H3. Il
\cref{ch:discussion} interpreta i risultati ed enuncia minacce alla
validità e limiti di scope; il \cref{ch:conclusions} conclude ed elenca
il lavoro futuro. Le appendici forniscono il registro condensato degli
esperimenti, le configurazioni complete, le tabelle di risultato per run
e i dettagli di riproducibilità.
